% Owner: member B. C supplies the final Transformer tuning values.
\paragraph{Sparse baseline.}
The sparse models use grouped five-fold cross-validation within the 5,909
training reviews. The TF-IDF grid covers minimum document frequency, a
vocabulary cap and optional character features. These control the trade-off
between coverage of rare expressions and the number of features. Logistic
regression uses $C=10$ and balanced class weights with 50,000 word unigram and
bigram features. The calibrated LinearSVC selects $C=1$ without class weights,
and ComplementNB selects $\alpha=0.1$ with 10,000 features.

\paragraph{Static embedding models.}
For mean GloVe, we compare $C\in\{0.1,1,10\}$ with and without balanced class
weights. At the common threshold of 0.5, the $C=10$ balanced setting has the
highest Macro-F1 (0.7623). Its unweighted counterpart scores 0.7598 Macro-F1
and 0.8988 accuracy, versus 0.8539 accuracy with weighting. The higher negative
F1 with weighting (0.6148 versus 0.5770) explains the selection despite the
accuracy trade-off. The BiLSTM begins with Adam, learning rate $10^{-3}$, hidden width 128
per direction, dropout 0.3, trainable embeddings, and ordinary cross-entropy.
A sequential search changes learning rate ($10^{-3}$ or $5\times10^{-4}$),
hidden width (64, 128 or 256), dropout (0.3 or 0.5), embedding freezing, then
class weighting. This limits cost and isolates each comparison near the
incumbent configuration. A factorial search would be needed to estimate
interactions between settings.
The learning rate controls step size, dropout reduces overfitting, and
freezing tests whether source supervision benefits from updating GloVe.
Class weighting changes gradients during training, while threshold tuning
changes the operating point after training. Each run uses batches of 64,
gradient clipping at 1.0, up to 10 epochs and patience of two epochs.

Candidates are ranked by validation Macro-F1 at threshold 0.5. The selected
BiLSTM configuration is repeated with seeds 42, 43 and 44. We use the
standard deviation to describe run-to-run variability. Since the same
validation fold also guides tuning and early stopping, these are model-selection
scores. The unlabeled product test set is used for prediction.
For each model, the accuracy and Macro-F1 thresholds are then fitted separately
on the validation probabilities.

The sequential sweep ran seven configurations, followed by the three final fits. The selected learning rate is $0.001$, with width 128 per direction, dropout 0.3, trainable embeddings, and ordinary cross-entropy. The single-seed Macro-F1 is $0.8383\pm0.0050$ (mean and population standard deviation). The three-seed probability average reaches 0.8524 Macro-F1 at threshold 0.5. Its validation-selected thresholds are $\tau_{\mathrm{acc}}=0.500$ and $\tau_{\mathrm{F1}}=0.500$. The individual settings and scores are recorded in \texttt{outputs/bilstm/hparam\_log.csv}.

\paragraph{Transformer.}
RoBERTa-base was fine-tuned with AdamW (learning rate $2\times10^{-5}$, weight decay 0.01), unweighted cross-entropy, batch size 4, maximum length 256 and a single epoch, in fp32 with a constant learning rate.
The small learning rate is the usual choice for fine-tuning a pre-trained encoder: larger steps overwrite the pre-trained representation, while one epoch was enough to reach about 95\% validation accuracy on 5.9k reviews.
The batch size of 4 was set by the memory of the laptop GPU.
Because of the deadline, no learning-rate or loss-weighting search was run for this model; class imbalance is handled by tuning the decision threshold on the validation split instead.
Seeds 42 and 43 reached validation Macro-F1@0.5 of 0.897 and 0.884. The seed-44 run diverged to a constant prediction (Macro-F1 0.461), a known failure mode of fine-tuning without learning-rate warm-up; a linear warm-up would be the first change to make.
\tbd{C: update if seed 44 is re-run with warm-up.}
